\chapter*{ABSTRAK}
\addcontentsline{toc}{chapter}{ABSTRAK}

\begin{center}
    \textbf{Pengembangan Sistem Manajemen Karyawan Terpusat Berbasis Integrasi Perangkat Kontrol Akses pada Bangunan Gedung Cerdas} \\
    \vspace{0.4cm}
    oleh \\
    \vspace{0.4cm}
    Axelius Davin\\
    18222016\\
\end{center}

\begin{spacing}{1.0}
Manajemen akses fisik merupakan elemen krusial dalam operasional bangunan gedung cerdas dengan mobilitas tinggi. Saat ini, sistem kontrol akses di Gedung ITB Innovation Park (IIP) masih bersifat terfragmentasi karena ketergantungan pada perangkat lunak bawaan vendor yang tertutup, sehingga menghambat visibilitas data pergerakan penghuni secara \textit{real-time}. Penelitian ini bertujuan untuk merancang dan mengimplementasikan sistem manajemen karyawan terpusat menggunakan arsitektur lapisan penengah (\textit{middleware}) untuk mengatasi hambatan interoperabilitas tersebut. Sistem dikembangkan menggunakan kerangka kerja NestJS, PostgreSQL, dan Next.js, dengan memanfaatkan protokol \textit{Intelligent Security API} (ISAPI) untuk menjembatani komunikasi antara peladen dengan terminal biometrik pengenal wajah. Metodologi yang digunakan adalah \textit{Design Thinking} yang mencakup tahapan investigasi kebutuhan hingga pengujian sistem. Hasil evaluasi menunjukkan bahwa sistem berhasil melakukan sinkronisasi data identitas tekstual secara terpusat dan menyajikan log riwayat akses secara \textit{real-time} dengan rata-rata latensi sebesar 11.285 ms (sekitar 11,28 detik). Selain itu, sistem mengimplementasikan mekanisme deteksi anomali yang secara otomatis menangkap bukti visual (\textit{snapshot}) saat terjadi pelanggaran akses, dengan tetap menjaga privasi data biometrik pada kondisi normal. Penelitian ini membuktikan bahwa pendekatan integrasi melalui protokol terbuka mampu meningkatkan efisiensi administratif dan visibilitas operasional dalam pengelolaan keamanan gedung cerdas.

Kata kunci: Gedung Cerdas, Kontrol Akses, \textit{Middleware}, ISAPI, \textit{Real-time Monitoring}.
\end{spacing}
